\section{Úvod}

Používateľské rozhranie predstavuje jednu z najdôležitejších častí softvérového systému, keďže priamo ovplyvňuje spôsob, akým používatelia so systémom interagujú. 
Aj technicky kvalitne navrhnutý systém môže zlyhať, pokiaľ jeho rozhranie nie je dostatočne intuitívne, prehľadné a zrozumiteľné.

V kontexte informačných systémov na správu dokumentov (DMS) je použiteľnosť používateľského rozhrania obzvlášť dôležitá. Tieto systémy sú často využívané rôznorodými skupinami používateľov, 
od bežných zamestnancov až po administrátorov, pričom každý z nich má odlišné požiadavky, skúsenosti a očakávania. Neintuitívne alebo neprehľadné rozhranie môže viesť k chybám, zníženiu produktivity alebo dokonca k odmietaniu systému používateľmi.

Cieľom tohto dokumentu je analyzovať používateľské rozhranie vybranej webovej aplikácie typu DMS prostredníctvom základného používateľského testovania. 
Testovanie je zamerané na identifikáciu problémových miest v rozhraní, overenie intuitívnosti jednotlivých krokov a získanie spätnej väzby od používateľov pri plnení konkrétnych úloh.

Testovanie je realizované na základe vopred definovaných scenárov použitia, ktoré reprezentujú typické operácie vykonávané v systéme, ako je správa číselníkov, vytváranie záznamov alebo spracovanie workflow úloh. 
Súčasťou testovania je aj formulácia testovacích úloh, pozorovanie správania používateľov a vyhodnotenie výsledkov.

Výsledkom dokumentu je súhrn identifikovaných pozitívnych a negatívnych aspektov používateľského rozhrania spolu s odporúčaniami na jeho zlepšenie. 
Tieto poznatky môžu slúžiť ako podklad pre ďalší rozvoj systému a zvýšenie jeho používateľskej kvality.

\clearpage
